\chapter{智能表示的双重基底：关系结构与属性内容}
\label{chp:dual_basis}

我们从一个常见而容易被忽略的现象开始：知道对象具有哪些属性，并不等于知道对象之间发生了什么；知道对象之间怎样关联，也不等于知道关系中承载着什么内容。一个系统可以准确识别杯子、桌面与水，却仍然无法判断杯子是否能够盛水。它也可以保存“容器包含内容”的规则，却无法在当前图像中找到需要操作的容器。前一种缺失关系及约束，后一种缺失属性及其绑定。智能的表示基础必须让二者能够共同进入预测、推理与行动。

我们把关系的组织方式称为“形”，把关系位置上承载的属性内容称为“质”。这两个名称延续形质构成论的表达，但在本书中首先指向计算对象。我们不把它们设定为宇宙的两个实体，也不把计算属性直接解释为主观体验。我们的任务是建立一种可以承载不同智能系统的表示语言，并说明这套语言怎样支持状态更新。它将成为 HSF-HD 状态动力学的基础，随后再由 AgentOS 落实为可维护的工作状态、历史记录与持久结构。

\section{从对象属性到关系结构：表示问题的起点}
\label{sec:physics_dilemma}

我们考虑两个包含相同对象的场景。第一个场景中，杯子位于桌面上，水位于杯子中；第二个场景中，杯子位于桌面旁，水流到桌面上。若系统只保存“杯子、桌子、水”三个类别及其颜色，两个场景可能得到相同的属性集合。然而，稳定放置、液体承载与下一步操作都发生了变化。由此可以看到，关系不是在属性列表之后添加的一层修辞，而是决定预测与行动是否有效的一类必要信息。我们需要保存对象身份、关系类型、关系方向以及关系成立的条件。

令 $V=\{1,\ldots,n\}$ 为当前表示中的位置集合。位置可以对应对象、事件、概念或被建模的局部单元，其含义由任务声明。令 $\mathcal{R}$ 为关系类型集合，$S^{(r)}\in\mathbb{R}^{n\times n}$ 表示类型 $r$ 的关系参数，$X\in\mathbb{R}^{n\times d}$ 表示各位置上的属性。我们据此定义表示对象
\begin{equation}
 \mathcal{Z}=\left(V,\{S^{(r)}\}_{r\in\mathcal{R}},X\right).
\end{equation}
这里的关系参数可以是离散关系的编码，也可以是带置信度的连续权重。它们是否对称、是否非负、是否具有概率意义，需要分别规定。若“包含”关系具有方向，就不能为了便于矩阵计算而把它强制改成无向关系。

杯子的颜色、容量估计与材质描述可以进入同一个对象的属性记录，但不同量不能在没有预处理的情况下直接相加。颜色类别、毫升数值与语义向量具有不同的类型和尺度。我们可以分别设置编码函数，把它们送入统一的计算空间，同时保留原始值、单位和来源。统一计算空间的作用是允许模块之间运算，不是宣布这些量在对象层面已经成为同一种属性。系统若要根据容量判断能否盛水，读出仍必须回到相应的数量及其不确定性。

位置索引也不是对象本身。系统内部把杯子编为位置一或位置三，不应改变任务结论。设 $P$ 是位置排列矩阵，则同一个表示在重新编号后满足 $X'=PX$、$S'^{(r)}=PS^{(r)}P^\top$。若任务输出不需要逐位置编号，我们要求读出函数 $g$ 满足
\begin{equation}
 g\left(\{PS^{(r)}P^\top\}_r,PX\right)
 =g\left(\{S^{(r)}\}_r,X\right).
\end{equation}
若输出需要指出操作对象，则应采用相应的等变关系，而不是把对象身份从输出中抹去。这个要求来自编号的任意性，其成立可以通过交换索引后的输出比较来检验。

关系与属性的分离并不要求它们统计独立。杯子的形态可能影响容量，事件的角色关系可能限制可填入的对象类型。我们的分解是在模型中明确二者的职责，使关系操作与属性操作具有可追踪的作用对象。当系统把“杯子可盛水”推广到另一个容器时，它可能保留容纳关系而改变材质与形状属性；当系统识别同一杯子被移到另一位置时，它可能保留对象属性而改变相对位置关系。两种变化需要分别表达，才能分析错误来自哪一部分。

我们由此得到一种有限而明确的建模立场：任务相关的表示应同时说明组织关系与局部内容。该立场并不意味着所有环境细节都必须被完整编码。对于寻找杯子的任务，关系可以较为粗略；对于安全搬运盛有热水的杯子，承载、接触、温度与动作约束就不能省略。我们通过任务需求决定表示的粒度，并把未被表示的因素纳入误差与边界讨论。形质分解首先帮助我们提出这些选择，而不是替代选择本身。

\section{结构与内容的绑定：从符号到分布式表示}
\label{sec:cognition_dilemma}

在语言中，绑定问题同样直接出现。“甲把钥匙交给乙”与“乙把钥匙交给甲”包含相同的对象词和动作词，但它们表达不同的事件。如果编码只对词向量求和，词序与角色可能消失；如果编码只保存“交付者、物品、接收者”的空槽位，却没有把实体内容写入槽位，事件仍然没有完成表示。我们需要让角色关系与属性内容共同决定表示，并在变换中保持对应。关系结构承担角色组织，内容记录承担实体描述，绑定机制使二者能够共同被读取。

我们用三个角色位置表示交付事件，分别对应交付者、被交付物与接收者。令 $S$ 描述角色之间的关系，令 $X_i$ 存储角色位置 $i$ 上的实体编码。若把甲与乙交换，改变的是内容与角色之间的对应，而不是仅仅重新编号。前一节的位置排列会同时变换结构和内容，保持同一事件；这里则在固定角色结构下交换两行内容，得到另一个事件。两种操作在代数上可以区分，这一点对于判断模型是否真正保存角色信息十分重要。

一个具体的联合编码可以写为 $e(o)=(S(o),X(o))$，其中 $o$ 是输入观测。读出函数给出预测 $\widehat y=g(S(o),X(o))$。我们可用任务损失训练联合编码，但仅凭预测误差较小，并不能唯一确定内部结构的解释。不同编码可能产生相同的任务输出。例如，一个模块可以把角色信息写入关系矩阵，另一个模块可以把角色信息混入内容向量。若我们需要对结构进行编辑，就必须增加编辑一致性、角色交换测试或对应的监督约束，使内部职责能够被观察和检验。

假设我们把关系类型与内容类型分别约束在给定集合中，并定义交换操作 $T$。若交换交付者和接收者的内容应使输出从“甲交付”变为“乙交付”，我们可以比较 $g(S,TX)$ 与任务所要求的变换结果。这个比较关注的是干预后的行为，而不是某一隐藏向量是否看起来像角色编码。对多对象场景，我们还需要检查同一种属性是否被错误地传播到其他对象。例如，红杯子放在蓝桌面上时，模型不能因为二者相邻，就把颜色一并写入桌面记录。

分布式表示并不天然排斥这种绑定。一个实体可以由多个位置上的活动共同编码，一个关系也可以由高维参数而不是单个符号保存。我们关心的是：在读出和编辑时，系统是否能够区分关系身份、内容身份及其对应。符号图、张量表示、神经网络与检索系统都可以实现这种要求，只是可观测接口不同。符号图较容易显示关系，连续向量较容易表达渐变相似性；二者的优势可以通过明确接口结合，而不需要宣布某一种实现已经在理论上取消另一种实现。

绑定还具有时间维度。系统在某一时刻识别到钥匙位于甲手中，随后观测到交付事件，就应更新钥匙的持有者，而不把旧状态与新状态混为同一事实。我们因此区分对象身份、事件身份与当前属性。历史记录保留变化的证据，工作状态维护当前任务所需的估计，持久结构保存交付关系的规律。在 AgentOS 中，这些对象可以分别进入情景记忆、工作记忆与结构记忆；这种工程映射需要根据具体实现确定，并不要求三种记忆采用同一数据结构。

我们采用形质分解，是为了建立可以诊断的联合表示。属性识别错误、关系识别错误、绑定错误和时间更新错误可以产生相似的最终失败，但它们需要不同的修复机制。一个系统若把错误归因只集中到输出层，就可能反复增加生成约束，却仍然无法修复内部的角色混淆。通过分别检查关系与内容，我们把问题从“表示是否足够强大”进一步转化为“哪些对象被保存、哪些对应被维护、哪些更新发生了错误”。

\section{从静态表示到状态更新：耦合与尺度}

形质分解在静态表示中区分关系与属性，在动力学中则进一步区分当前活动和持久组织。我们让 $X_k$ 表示第 $k$ 个更新时刻的活动状态，$S_k$ 表示当时可用的关系结构。最一般的离散模型可以写为 $X_{k+1}=F(X_k,S_k,u_k,c_k)$ 与 $S_{k+1}=G(S_k,X_k,m_k,v_k)$，其中输入 $u_k$ 提供外部信息，控制 $c_k$ 调节当前更新，历史变量 $m_k$ 提供已有经历，目标变量 $v_k$ 规定任务偏好。这里每个函数的输入和输出空间必须明确，不能仅凭符号相邻就把不同对象相乘。

为了分析一种具体机制，我们暂时固定位置数量与属性维度，并假设 $S$ 是对称非负权图的邻接矩阵。令 $D_{ii}=\sum_j S_{ij}$，$L_S=D-S$。对任何 $q\in\mathbb{R}^n$，利用对称性可得
\begin{align}
 q^\top L_Sq
 &=\sum_i\Bigl(\sum_j S_{ij}\Bigr)q_i^2-\sum_{i,j}S_{ij}q_iq_j\\
 &=\frac12\sum_{i,j}S_{ij}(q_i-q_j)^2\geq0.
\end{align}
因此，$L_S$ 是半正定矩阵。这个性质使我们可以把 $-L_SX$ 用作当前结构上的差异消减机制。若关系具有方向或带负权，以上证明的条件就不成立，需要改用相应算子并重新分析。

现在考虑固定结构下的连续近似 $\dot X=-L_SX-\gamma X+Bu+Uc$，其中 $\gamma>0$。在这个矩阵状态模型中，我们规定 $u\in\mathbb{R}^{d_u\times d}$、$c\in\mathbb{R}^{d_c\times d}$、$B\in\mathbb{R}^{n\times d_u}$ 和 $U\in\mathbb{R}^{n\times d_c}$，使每个驱动项都属于 $\mathbb{R}^{n\times d}$。取 $V(X)=\tfrac12\|X\|_F^2$，在输入与控制为零时有
\begin{equation}
 \dot V=-\operatorname{tr}(X^\top L_SX)-\gamma\|X\|_F^2
 \leq-2\gamma V.
\end{equation}
由微分不等式可得 $V(t)\leq e^{-2\gamma t}V(0)$。我们在这里证明的是一个指定线性系统在零输入条件下的活动衰减，不能把它解释为系统理解了任务，也不能由衰减推出正确决策。事实上，所有活动趋于零可能使任务完全停止。

如果输入项持续存在，系统的轨迹由输入驱动，零状态通常不再是平衡状态。我们需要研究输入大小、扰动范围与响应之间的关系。若控制又依赖当前状态，例如 $c=KX$，闭环矩阵变为 $-L_S-\gamma I+UK$。原来的衰减结论不能直接沿用，因为反馈项可能抵消甚至超过阻尼。我们必须重新检查闭环算子的对称部分或寻找适当的正定度量。这一例子说明，增加一个被称为“控制”的模块，并不自动使系统更加稳定。

为了把连续机制落实为离散计算，我们还需要选择更新步长。对于无输入的线性过程，令 $A=L_S+\gamma I$，显式更新为 $X_{k+1}=(I-hA)X_k$。由于 $A$ 对称正定，逐特征方向检查得到条件 $|1-h\lambda_i(A)|<1$，因此要求 $0<h<2/\lambda_{\max}(A)$。连续轨迹能够衰减，并不意味着任意步长的程序都能够衰减。步长过大时，离散轨迹可以交替放大；我们必须把积分器的条件纳入实现，而不能只在理论说明中写出一个稳定的微分方程。

在持续文档处理的场景中，同样的问题表现为更新频率与反馈延迟的配合。一个控制器根据旧状态提高某类证据的优先级，而新的证据已经改变了当前结论，这次控制就可能强化过时的活动。我们可以保存控制所依据的状态版本，并在执行之前核对其有效性。版本检查不能替代所有延迟分析，但它使“控制依据是什么”成为可以追踪的事实，从而避免把异步系统误当成同步更新系统。

结构学习使问题进一步变化。我们可以让关系权重按任务损失更新，但当前活动、记忆与目标都在变化时，损失沿完整轨迹的导数包含所有分量。结构方向上的梯度下降只能保证该方向的局部下降，不能单独保证整个耦合系统的目标单调下降。若我们引入较小的结构更新率 $\varepsilon_S$，它表示一种速度选择；只有当快速过程能够在结构显著变化之前充分响应，尺度分离近似才具有合理依据。速度之比、延迟与扰动都需要被测量。

我们把这些限制视为动力学理论的组成部分。一个状态方程需要说明它承载什么活动，一个结构方程需要说明它改变什么组织，一项稳定性结论需要说明它针对哪些固定或变化的变量。在后续章节中，我们会分别处理历史保持、目标控制与跨尺度耦合。本章不提前给出所有智能系统的统一方程，而是建立一种推导纪律：先声明对象和条件，再给出更新机制，最后判断这个机制实际能够支持哪些结论。

\section{模型的解释边界与可验证要求}
\label{sec:axiom_system}

我们把本章的论述整理为一组可以使用、也可以被修正的建模约定。首先，任务相关的对象可以通过关系结构与属性内容共同表示；其次，关系与属性的绑定必须在读取、编辑和更新中保持可追踪性；再次，结构与活动可以采用不同的更新规则，并通过明确接口耦合。这些约定为模型构造提供起点，但它们不等于已经得到经验确认的普遍规律。是否适用于一个具体智能系统，需要进一步分析其观测、任务与实现条件。

“形”并不要求一定采用连续流形。在有限知识图中，它可以由离散关系组成；在空间感知中，它可以包含坐标、接触与遮挡；在程序执行中，它可以包含调用关系与依赖顺序。我们只在需要连续近似、局部坐标或度量时，引入相应几何结构。同样，“质”也不要求一定采用复数场。类别标记、数值属性、概率分布、连续特征和多模态编码都可以成为内容记录。采用哪一种对象，要由任务与计算机制说明，而不是由名称决定。

在解释一个表示时，我们还需要区分可预测性与可辨识性。若模型能够准确预测杯子位置，说明它在指定数据与任务上获得了有效信息，却不能证明它内部的每一维都对应真实空间中的某一属性。若模型能够通过结构编辑改变预测，我们才得到关于编辑接口的进一步证据。即使如此，编辑接口的有效性也可以是局部的：在训练分布内可用，不代表在任意对象组合下都成立。因此，模型评估必须说明分布、干预范围和失败条件。

我们可以设置四类互补的检查。对象检查考察身份是否持续、属性是否保持类型与单位；关系检查考察方向、组合与约束是否正确；绑定检查考察角色交换、相邻干扰与来源追踪；更新检查考察观测变化后状态是否一致，以及旧信息是否被保留为历史而不是误当成当前事实。这些检查不需要每次都采用复杂指标，但必须让失败具有可解释的归属。只有最终回答正确率而没有中间检查，往往难以区分偶然成功与机制成立。

对于杯子场景，系统可以先建立对象记录，再建立“位于”“包含”和“接触”关系，随后根据新观测更新位置与液体状态。我们可以故意交换颜色、遮挡杯把、改变桌面倾斜或延迟传感输入，观察模型是否仍能保持正确的对象绑定。这里的干预不是为了制造一个困难示例，而是为了区分表面相关性与任务所需的关系。若颜色改变就使容量判断失效，说明系统可能依赖了不应决定容量的特征；若来源信息被丢失，则需要修复记录与绑定机制。

理论层与工程层由此具有不同职责。HSF-HD 说明哪些对象共同构成状态、哪些更新发生耦合以及可以在什么条件下得到性质。AgentOS 负责维护记录、执行更新、监控资源和组织交互。一个工程实现可以把结构分布在模型参数、索引与规则中，把活动分布在工作记忆和模块缓存中；它需要提供从实现对象到理论变量的读出方式。我们不要求二者在存储位置上一一对应，但要求理论中的区分能够在工程中找到可观测的实现含义。

本章建立的双重基底，最终服务于后续的状态动力学，而不是在表示层结束讨论。结构规定当前有哪些可用联系，属性决定联系中承载哪些内容，绑定保持它们的共同身份，更新机制决定这种联合表示怎样随输入与目标变化。我们将在下一章专门讨论关系结构的构造与操作，再讨论属性状态的尺度、激活与读出。这样，智能的持续生成可以从清楚的计算对象开始，而不需要借助未经限定的物理等同来补足论证。
